\section{Method: BRADD}
\label{sec:method}

\begin{figure*}[h]
    \centering
    \includegraphics[width=0.7\textwidth]{Images/BRIDGE.pdf}
    \captionsetup[subfigure]{skip=0pt}
    \caption{The BRADD algorithm augments the most underrepresented datapoints in a dataset based on OOD detection.}
    \label{fig:full_pipeline}
\end{figure*}

We propose BRADD (Balancing Representations with Anomaly Detection and Diffusion), a data-centric approach to address imbalance in self-supervised learning. Unlike model-centric approaches, BRADD identifies and augments underrepresented concepts in the data distribution itself.

\begin{algorithm}[h]
\caption{BRADD: Balancing Representations through Anomaly Detection and Diffusion}
\label{alg:bradd}
\begin{algorithmic}[1]
\STATE {\bfseries Input:} model $m$, SSL algorithm $a$, diffusion model $m_d$, cycles $N_C$, epochs per cycle $N_E$, dataset $D$, samples per point $N_{Aug}$, OOD samples $N_{OOD}$
\REPEAT
\STATE Train $m$ using $a$ for $N_E$ epochs on $D$
\STATE Compute embeddings $E = \{m(x_i) | x_i \in D\}$
\STATE Find OOD points $O = \{x_i | \text{rank}(d_{kNN}(m(x_i), E)) \leq N_{OOD}\}$
\FOR{$x \in O$}
\STATE $D = D \cup \{m_d(x, \epsilon_j) | j = 1,...,N_{Aug}\}$
\ENDFOR
\UNTIL{$N_C$ cycles completed}
\end{algorithmic}
\end{algorithm}

\noindent
BRADD alternates between self-supervised pre-training and dataset augmentation phases (Algorithm \ref{alg:bradd}). After training for $N_E$ epochs, we identify underrepresented data points by computing k-nearest neighbor distances in the embedding space. Rather than using a percentile threshold, we select the top-$N_{OOD}$ samples with highest k-NN distances, providing precise control over augmentation.

\noindent
For each identified OOD point, we generate $N_{Aug}$ new samples using Stable Diffusion 2 UnCLIP \cite{rombach2021highresolution}, which preserves semantic content while introducing sufficient variation. We use k=5 for k-NN computation, $N_{OOD}=500$, and $N_{Aug}=5$, adding 2,500 new images per cycle across 5 cycles with 20 epochs each (100 total epochs).